% Problem record op_09dc815c8a194fb7. This TeX file is derived from
% database/problems_json/op_09dc815c8a194fb7.json by scripts/sync-tex.mjs; edit the JSON record
% and rerun the script rather than editing this file.
\section{Distillable entanglement of the antisymmetric state}
\label{sec:op_09dc815c8a194fb7}

\paragraph{Problem.}
What is the exact LOCC distillable entanglement
$E_D(\rho_{\mathrm{as}}^{(d)})$ of the $d\times d$ antisymmetric state, for
each integer $d\ge3$ --- as a closed form, or asymptotically through the
existence and value of $\lim_{d\to\infty}d\,E_D(\rho_{\mathrm{as}}^{(d)})$?

Let $F$ be the swap operator on $\mathbb C^d\otimes\mathbb C^d$, acting by
$F\,|ij\rangle=|ji\rangle$ on the computational basis.  Define the
antisymmetric projector and the normalized antisymmetric state
\begin{equation}
  P_{\mathrm{as}}:=\tfrac12(\mathbb 1-F),
  \qquad
  \rho_{\mathrm{as}}^{(d)}:=\frac{2}{d(d-1)}\,P_{\mathrm{as}},
  \label{eq:09dc-antisymmetric-state}
\end{equation}
so Eq.~\eqref{eq:09dc-antisymmetric-state} is the extreme point $\alpha=-1$
of the Werner family $\rho_\alpha\propto\mathbb 1+\alpha F$ and is not of
positive partial transpose; for $d=2$ it is the singlet, with $E_D=1$, and
the question concerns $d\ge3$.

Let $\Phi_2$ be a maximally entangled state of Schmidt rank two.  The LOCC
distillable entanglement of the state in
Eq.~\eqref{eq:09dc-antisymmetric-state} is
\begin{equation}
  E_D\bigl(\rho_{\mathrm{as}}^{(d)}\bigr):=\sup\Bigl\{r\ge0:\ \exists\ 
  \mathrm{LOCC}\ \Lambda_n\ \mathrm{with}\ 
  \lim_{n\to\infty}\bigl\|\Lambda_n\bigl((\rho_{\mathrm{as}}^{(d)})^{\otimes n}\bigr)
  -\Phi_2^{\otimes\lfloor rn\rfloor}\bigr\|_1=0\Bigr\},
  \label{eq:09dc-distillable-entanglement}
\end{equation}
and replacing LOCC in Eq.~\eqref{eq:09dc-distillable-entanglement} by
completely PPT-preserving operations defines the relaxation
$E_{D,\mathrm{PPT}}(\rho_{\mathrm{as}}^{(d)})$.  Determine
$E_D(\rho_{\mathrm{as}}^{(d)})$ or $E_{D,\mathrm{PPT}}(\rho_{\mathrm{as}}^{(d)})$
exactly for any single $d\ge3$, or determine whether
$\lim_{d\to\infty}d\,E_D(\rho_{\mathrm{as}}^{(d)})$ exists and what its value
is.  A complete answer consists of an explicit achievable protocol together
with a matching converse bound in the chosen operation class.

\subsection*{Status}

Unsolved

\subsection*{Source}

Implicit in Christandl, Schuch, and Winter's analysis of the antisymmetric
state: they bound its distillable key, and hence its distillable entanglement,
by $O(1/d)$ while its entanglement cost stays above a $d$-independent
constant, leaving the exact rate open
\sourcecite{ref:09dc-christandl-schuch-winter}{CSW12}; announced in
\sourcecite{ref:09dc-csw-prl}{CSW10}.  Yura's exact evaluation of the
entanglement cost at $d=3$ settles the cost side of the same question
\sourcecite{ref:09dc-yura}{Yur03}.

\subsection*{Progress}

\begin{itemize}
  \item DiVincenzo, Shor, Smolin, Terhal, and Thapliyal determined the
  one-copy-distillable region of the Werner family: a Schmidt-rank-two
  projection witnesses entanglement in a single copy precisely on one side of
  an explicit threshold, and the state in
  Eq.~\eqref{eq:09dc-antisymmetric-state} lies in that region, so
  $E_D(\rho_{\mathrm{as}}^{(d)})>0$
  \sourcecite{ref:09dc-divincenzo-et-al}{DST+00}.

  \item The logarithmic negativity upper-bounds distillable entanglement
  \sourcecite{ref:09dc-vidal-werner}{VW02}, and the Rains bound upper-bounds
  distillation by completely PPT-preserving operations
  \sourcecite{ref:09dc-rains}{Rai01}.  The partial transpose of the state in
  Eq.~\eqref{eq:09dc-antisymmetric-state} has the single negative eigenvalue
  $-1/d$ and the eigenvalue $1/(d(d-1))$ with multiplicity $d^2-1$,
  \begin{equation}
    \operatorname{spec}\bigl(\bigl(\rho_{\mathrm{as}}^{(d)}\bigr)^{\Gamma}\bigr)
    =\bigl\{-\tfrac1d\bigr\}\cup\bigl\{\tfrac{1}{d(d-1)}\ \text{with
    multiplicity}\ d^2-1\bigr\},
    \label{eq:09dc-pt-spectrum}
  \end{equation}
  so its trace norm is $(d+2)/d$ and the converse bounds bracket the rates as
  \begin{equation}
    0<E_D\bigl(\rho_{\mathrm{as}}^{(d)}\bigr)
    \le E_{D,\mathrm{PPT}}\bigl(\rho_{\mathrm{as}}^{(d)}\bigr)
    \le\log_2\bigl(1+\tfrac2d\bigr).
    \label{eq:09dc-bracket}
  \end{equation}
  The upper bound in Eq.~\eqref{eq:09dc-bracket} follows from
  Eq.~\eqref{eq:09dc-pt-spectrum} and tends to zero as $d\to\infty$; whether
  it is attained is open.  The Rains bound itself --- and the regularized
  PPT relative entropy of entanglement --- is exactly known for this
  state: Audenaert, Eisert, Jan{\'e}, Plenio, Virmani, and De Moor showed that
  $E_{R,\mathrm{PPT}}^{\infty}(\rho_{\mathrm{as}}^{(d)})=\log_2\bigl(1+\tfrac2d\bigr)$
  \sourcecite{ref:09dc-audenaert-et-al}{AEA+01}, and the chain
  $E_{R,\mathrm{PPT}}^{\infty}=E_{\mathrm{Rains}}=E_N=\log_2\bigl(1+\tfrac2d\bigr)$,
  with $E_{\mathrm{Rains}}$ the single-copy Rains bound, is recorded in
  Sec.~III of \sourcecite{ref:09dc-christandl-schuch-winter}{CSW12}.

  \item On the cost side, Yura computed $E_C(\rho_{\mathrm{as}}^{(3)})=1$ ebit
  exactly \sourcecite{ref:09dc-yura}{Yur03}.  Christandl, Schuch, and Winter
  then showed that the distillable key obeys
  $K_D(\rho_{\mathrm{as}}^{(d)})=O(1/d)$, hence
  $E_D(\rho_{\mathrm{as}}^{(d)})\le K_D(\rho_{\mathrm{as}}^{(d)})\to0$ as
  $d\to\infty$, while the entanglement cost stays above a $d$-independent
  constant, making asymptotic manipulation of the family irreversible in
  the limit $d\to\infty$ --- the cited bounds separate $E_D$ from $E_C$
  only for $d\ge7$ (and at $d=3$, via Yura's $E_C=1$), leaving $d=4$,
  $5$, and $6$ open.  The best known lower bound is
  $E_D(\rho_{\mathrm{as}}^{(d)})\ge1/d$, recorded for both $E_D$ and $K_D$,
  making the $O(1/d)$ upper bound optimal up to a constant; the exact
  distillable entanglement is left open
  \sourcecite{ref:09dc-christandl-schuch-winter}{CSW12}.

  \item For any rank-two mixed state supported on the $3\otimes3$ antisymmetric
  subspace, Wang and Duan showed that the distillable entanglement under
  completely PPT-preserving operations is strictly smaller than one ebit,
  while the PPT entanglement cost is exactly one ebit, a strict separation
  between the two PPT measures; the inequality does not evaluate the rate in
  Eq.~\eqref{eq:09dc-distillable-entanglement}
  \sourcecite{ref:09dc-wang-duan}{WD17}.

  \item Exact rates are known only beyond the PPT class or beyond quantum
  operations: Lami and Regula computed the exact distillation rate under
  dually non-entangling operations, a class strictly containing PPT-preserving
  operations, so it does not determine either rate in
  Eq.~\eqref{eq:09dc-bracket} \sourcecite{ref:09dc-lami-regula}{LR24}; and
  Chen, Wang, Zhang, and Zhu established reversible exact state transfer under
  trace-preserving PPT-preserving maps that need not be completely positive,
  governed by logarithmic negativity, a zero-error non-CPTP setting distinct
  from the vanishing-error rates posed here
  \sourcecite{ref:09dc-chen-et-al}{CWZ+25}.

  \item Kondra, Brinster, Kampermann, Bru{\ss}, and Wyderka constructed
  analytically solvable families built from antisymmetric states exhibiting
  exponential strong-converse irreversibility under completely PPT-preserving
  operations \sourcecite{ref:09dc-kondra-et-al}{KBK+26}.  The results concern
  cost and error exponents rather than the vanishing-error rate in
  Eq.~\eqref{eq:09dc-distillable-entanglement}, and the LOCC-side analogues
  are stated there as open.

  \item Sen and Lami computed the main entanglement measures of the flower-state
  family, whose distillable entanglement is exactly one ebit independent of
  local dimension, citing the antisymmetric state as the standard example of
  vanishing distillable entanglement at nonvanishing cost; the value for the
  antisymmetric state itself is not computed
  \sourcecite{ref:09dc-sen-lami}{SL26}.
\end{itemize}

\subsection*{References}

\begin{enumerate}[label={},leftmargin=4.8em,labelsep=0.5em]
  \footnotesize
  \item[\textup{[DST+00]}]\label{ref:09dc-divincenzo-et-al}
  D. P. DiVincenzo, P. W. Shor, J. A. Smolin, B. M. Terhal, and A. V.
  Thapliyal, ``Evidence for Bound Entangled States with Negative Partial
  Transpose,'' \emph{Physical Review A} \textbf{61}, 062312 (2000).
  \href{https://arxiv.org/abs/quant-ph/9910026}{arXiv:quant-ph/9910026}.

  \item[\textup{[AEA+01]}]\label{ref:09dc-audenaert-et-al}
  K. Audenaert, J. Eisert, E. Jan{\'e}, M. B. Plenio, S. Virmani, and
  B. De Moor, ``The Asymptotic Relative Entropy of Entanglement,''
  \emph{Physical Review Letters} \textbf{87}, 217902 (2001).
  \href{https://doi.org/10.1103/PhysRevLett.87.217902}{doi:10.1103/PhysRevLett.87.217902};
  \href{https://arxiv.org/abs/quant-ph/0103096}{arXiv:quant-ph/0103096}.

  \item[\textup{[Rai01]}]\label{ref:09dc-rains}
  E. M. Rains, ``A Semidefinite Program for Distillable Entanglement,''
  \emph{IEEE Transactions on Information Theory} \textbf{47}, 2921--2933
  (2001).
  \href{https://arxiv.org/abs/quant-ph/0008047}{arXiv:quant-ph/0008047}.

  \item[\textup{[VW02]}]\label{ref:09dc-vidal-werner}
  G. Vidal and R. F. Werner, ``Computable Measure of Entanglement,''
  \emph{Physical Review A} \textbf{65}, 032314 (2002).
  \href{https://arxiv.org/abs/quant-ph/0102117}{arXiv:quant-ph/0102117}.

  \item[\textup{[Yur03]}]\label{ref:09dc-yura}
  F. Yura, ``Entanglement Cost of Three-Level Antisymmetric States,''
  \emph{Journal of Physics A: Mathematical and General} \textbf{36}, L237
  (2003).
  \href{https://doi.org/10.1088/0305-4470/36/15/104}{doi:10.1088/0305-4470/36/15/104};
  \href{https://arxiv.org/abs/quant-ph/0302163}{arXiv:quant-ph/0302163}.

  \item[\textup{[CSW10]}]\label{ref:09dc-csw-prl}
  M. Christandl, N. Schuch, and A. Winter, ``Highly Entangled States With
  Almost No Secrecy,'' \emph{Physical Review Letters} \textbf{104}, 240405
  (2010).
  \href{https://doi.org/10.1103/PhysRevLett.104.240405}{doi:10.1103/PhysRevLett.104.240405}.

  \item[\textup{[CSW12]}]\label{ref:09dc-christandl-schuch-winter}
  M. Christandl, N. Schuch, and A. Winter, ``Entanglement of the
  Antisymmetric State,'' \emph{Communications in Mathematical Physics}
  \textbf{311}, 397--422 (2012).
  \href{https://doi.org/10.1007/s00220-012-1446-7}{doi:10.1007/s00220-012-1446-7};
  \href{https://arxiv.org/abs/0910.4151}{arXiv:0910.4151}.

  \item[\textup{[WD17]}]\label{ref:09dc-wang-duan}
  X. Wang and R. Duan, ``Irreversibility of Asymptotic Entanglement
  Manipulation Under Quantum Operations Completely Preserving Positivity of
  Partial Transpose,'' \emph{Physical Review Letters} \textbf{119}, 180506
  (2017).
  \href{https://doi.org/10.1103/PhysRevLett.119.180506}{doi:10.1103/PhysRevLett.119.180506}.

  \item[\textup{[LR24]}]\label{ref:09dc-lami-regula}
  L. Lami and B. Regula, ``Distillable Entanglement under Dually
  Non-Entangling Operations,'' \emph{Nature Communications} \textbf{15},
  10120 (2024).
  \href{https://doi.org/10.1038/s41467-024-54201-5}{doi:10.1038/s41467-024-54201-5};
  \href{https://arxiv.org/abs/2307.11008}{arXiv:2307.11008}.

  \item[\textup{[CWZ+25]}]\label{ref:09dc-chen-et-al}
  Yu-Ao Chen, Xin Wang, Lei Zhang, and Chenghong Zhu, ``Reversible
  Entanglement beyond Quantum Operations,'' \emph{Physical Review Research}
  \textbf{7}, 013297 (2025).
  \href{https://doi.org/10.1103/PhysRevResearch.7.013297}{doi:10.1103/PhysRevResearch.7.013297}.

  \item[\textup{[KBK+26]}]\label{ref:09dc-kondra-et-al}
  T. V. Kondra, R. Brinster, H. Kampermann, D. Bru{\ss}, and N. Wyderka,
  ``Very Strong Irreversibility of Quantum Entanglement,'' arXiv preprint
  (2026).
  \href{https://arxiv.org/abs/2607.27195}{arXiv:2607.27195}.

  \item[\textup{[SL26]}]\label{ref:09dc-sen-lami}
  S. Sen and L. Lami, ``Entanglement of Flower States,'' arXiv preprint
  (2026).
  \href{https://arxiv.org/abs/2608.02587}{arXiv:2608.02587}.
\end{enumerate}

\subsection*{Comment}

Even at $d=3$ the exact value is unknown: Eq.~\eqref{eq:09dc-bracket}
leaves $0<E_D\le E_{D,\mathrm{PPT}}\le\log_2(5/3)$, while the LOCC
entanglement cost equals one ebit, so both the rate and the sharpness of the
irreversibility are undetermined.  Asymptotically, the sharp question is the
leading constant: $E_D(\rho_{\mathrm{as}}^{(d)})=O(1/d)$ by the key bound,
while the achievability side is known only up to constants
(\sourcecite{ref:09dc-christandl-schuch-winter}{CSW12} records a $1/d$ lower
bound on $E_D$, matching the $\log_2\bigl(1+\tfrac2d\bigr)\sim 2/(d\ln 2)$
converse up to a constant), so neither
the existence nor the value of
$\lim_{d\to\infty}d\,E_D(\rho_{\mathrm{as}}^{(d)})$ is established.  A proof that the LOCC and PPT rates differ on this family, or a
determination whether the known Rains value $\log_2\bigl(1+\tfrac2d\bigr)$ is
asymptotically attained by completely PPT-preserving distillation, would
settle one of the posed alternatives.
This record asks for the value of a rate on one explicit family, unlike
Problem~\ref{sec:op_b93eb7197c926d9e}, the LOCC distillable entanglement of
Bell-diagonal qubit states, and
Problem~\ref{sec:op_4cf3e7b8663b1d41}, whether the Rains bound of
Bell-diagonal states is attained by a PPT-preserving channel;
Problem~\ref{sec:op_75b91a20dd384110} asks for a closed form of the
regularized exact zero-error PPT distillable entanglement of an arbitrary
support, a different quantity from the vanishing-error rate in
Eq.~\eqref{eq:09dc-distillable-entanglement}.

\subsection*{Field}

Quantum Resource Theory

\subsection*{Topic}

Entanglement distillation; Entanglement measures; Local operations and classical communication; PPT-preserving operations

\subsection*{ID}

\texttt{op\_09dc815c8a194fb7}
